\subsection{Fixed row coverage and exposure-weighted coverage}
% Show the moved Proposition 1 without disturbing subsequent numbering.
\begingroup
\renewcommand{\theproposition}{1}
\renewcommand{\theHproposition}{threshold}
\begin{proposition}[Fixed-coverage threshold principle]
\label{prop:threshold}
Assume $\E|L-rW|<\infty$. Among rules with $\E a=c$, accepting
the smallest $m_r=\E[L-rW\mid X]$ values minimizes expected excess.
Boundary randomization resolves atoms. The minimum excess is
$\int_0^cQ_m(u)\,du$, where $Q_m$ is the quantile function of $m_r$.
A ratio-risk assertion additionally requires positive accepted exposure.
\end{proposition}
\endgroup
\addtocounter{proposition}{-1}
Choose $t$ with $\Pr(m_r<t)\leq c\leq\Pr(m_r\leq t)$, and let
$a^*$ accept below $t$, reject above it, and randomize on the tie
to attain $c$. For any other rule of the same coverage,
\[
 (a-a^*)(m_r-t)\geq0,\qquad
 \E[a(L-rW)]-\E[a^*(L-rW)]
 =\E[(a-a^*)(m_r-t)]\geq0.
\]
This proves the threshold rule and the lower-quantile expression.
The limiting rules cover $c=0,1$. This minimizes excess, not the
ratio itself at fixed row coverage, because accepted exposure can vary.

For Proposition~\ref{prop:weighted}, $w=0$ implies $W=0$ almost
surely conditionally. These contexts add no exposure: reject them
if $\ell>0$; otherwise they are neutral. On $\{w>0\}$, set
$d\nu=w\,dP_X/T$ and $\xi=\ell/w$. Then
\[
 \E[aW]=T\E_\nu a,\qquad \E[aL]=T\E_\nu[a\xi].
\]
The exchange argument under $\nu$ minimizes the numerator at fixed
exposure coverage by ordering $\xi$. Maximizing feasible exposure
coverage thus admits an $\ell/w$ threshold when a positive-exposure
feasible rule exists. An additional binding row floor can change
that optimizer. Randomization handles atoms; the ranking is
independent of $r$, while the selected threshold depends on $r$.
